\documentclass{article}
\usepackage[utf8]{inputenc}
\usepackage{amsmath}
\usepackage{amssymb}
\usepackage{booktabs}
\usepackage{geometry}
\geometry{a4paper, margin=1in}
\usepackage{hyperref}
\usepackage{comment}
\hypersetup{
    colorlinks=true,
    linkcolor=blue,
    filecolor=magenta,
    urlcolor=cyan,
}
\usepackage{url}
\usepackage{algorithm}
\usepackage{algpseudocode}
\usepackage{listings}
\lstset{
    showstringspaces=false,
}
\usepackage{float}
\usepackage{makecell}
\renewcommand\theadfont{\bfseries}
\usepackage[dvipsnames]{xcolor}
\usepackage{tikz}
\usetikzlibrary{shapes.geometric, arrows.meta,
  positioning}

\DeclareMathOperator{\median}{median}

% Defines a command for author comments.
% Usage: \authcomment{This is my comment.}
\newcommand{\authcomment}[1]{\textcolor{Red}{\textbf{[#1]}}}

% Defines an environment for author comments that can contain listings and other complex content
% Usage: \begin{authcommentenv} ... \end{authcommentenv}
\newenvironment{authcommentenv}
  {\par\begingroup\color{Red}\bfseries[}% Begin code: red bold and opening bracket
  {]\endgroup\par}% End code: closing bracket and end group

\title{Final CTL Instruction on Training PINNs for the Three-Body Problem}
\author{Starlordpd233}
\date{June 2025}
